% ---------------------------------------------------------------
\section{Journal--institution citation networks}
\label{sec:networks}
% ---------------------------------------------------------------

%\paragraph{Corpus construction.}
We begin corpus construction by selecting $N_s$ sources $\mathfrak{S}$ in a field, discipline or category. 
Then we select $N_w$ works $\mathfrak{W}$ published in $\mathfrak{S}$ within the census 
window years $[t^c_0,t^c_1]$ or the earlier contiguous target year $t^t$ 
where $t^t = t^c_0 - 1$. 
From the reference lists of the set of census works $\mathfrak{W}^c$ we identify the set of target works $\mathfrak{W}^t$ and form the filtered reference set  $\mathfrak{R}=\{(i,j):w_i\in\mathfrak{W}^c,\,w_j\in\mathfrak{W}^t\}$.
Author affiliations in $\mathfrak{W}$ establish a set of institutions $\mathfrak{U}'$ and a set of 
authors $\mathfrak{A}'$. Examining the frequency distributions of institutions and authors, we retain institutions with work counts $\geq\tau^U$ and authors with work counts $\geq\tau^A$, giving $\mathfrak{U}$ and $\mathfrak{A}$ respectively. 
Each work $w_i$ has an \textit{institution weight} under author-fractional counting,
\begin{equation}
  \omega_{ik} \;=\;
  \sum_{j:\,k\in\mathcal{u}_{ij}}
  \frac{1}{a_i}\,\frac{1}{u_{ij}},
  \label{eq:omega}
\end{equation}
where $a_i$ is the number of authors and $u_{ij}$ the number of
institutions of author $j$. Since $\omega_{iu}$ requires both an author and an institution, we drop any authorship where either the author $\notin\mathfrak{A}$ or the institution $\notin\mathfrak{U}$, and subsequently drop works that lose all their authorships. The institution weights are not renormalised to account for the omission of an authorship, so that $\sum_k\omega_{ik}\leq 1$.

When a work $w_i$ references another work $w_j$ we say that $w_i$ gives attention to $w_j$. The attention given by a work can be either (a) the number of references, or (b) a fixed amount \citep{zittModifyingJournalImpact2008}.  Let $R_i=|\{j:(i,j)\in\mathfrak{R}\}|$ be the reference count of work $w_i$ and $\bar{R}=|\mathfrak{R}|/N_w$ the mean reference count over $\mathfrak{W}$. We use the \textit{reference normalisation weight} $\delta_i$ to select case (a) or case (b) respectively:
\begin{equation}
	\delta_i \;=\; 
	\begin{cases}
		 1 & \text{(full count)} \\ 
		\bar{R}/R_i & \text{(fixed count).} 
	\end{cases}
	\label{eq:delta}
\end{equation}

We construct the set of \emph{units} [sources or institutions;  \cite{pinskiCitationInfluenceJournal1976}]
$\mathfrak{P}=\mathfrak{S}\cup\mathfrak{U}$ with $N=N_s+N_u$ elements, ordered so that 
sources occupy indices $1,\dots,N_s$ and institutions indices $N_s+1,\dots,N$. 
A reference $(i,j) \in \mathfrak{R}$ gives attention across unit-type pairs: 
source--source ($SS$), source--institution ($SI$), institution--source ($IS$), and 
institution--institution ($II$).

For each work $w_i$ we define the source membership one-hot vector
$\mathbf{b}^S_i\in\mathbb{R}^{N_s}$ with $[\mathbf{b}^S_i]_k=\mathbf{1}[w_i\in s_k]$,
and the institution weight distribution vector
$\mathbf{b}^U_i\in\mathbb{R}^{N_u}$ with $[\mathbf{b}^U_i]_u=\omega_{iu}$
from \eqref{eq:omega}.
Attention flow between pairs of unit types is accumulated as four blocks over $\mathfrak{R}$ as rank-1 outer products,
\begin{equation}
	\mathbf{C}_{XY} \;=\; \sum_{(i,j)\,\in\,\mathfrak{R}}
	\delta_i\;\mathbf{b}^X_i\,\bigl(\mathbf{b}^Y_j\bigr)^\top,
	\qquad XY\in\{SS,\,SI,\,IS,\,II\},
	\label{eq:Cblock}
\end{equation}
with $\mathbf{b}^X_i$ weighting the referencing (citing) side and $\mathbf{b}^Y_j$ the referenced (cited) side.

The citation matrix assembles the four blocks:
\begin{equation}
	\mathbf{C}(\chi,\mathbf{m},\delta) \;=\;
	\begin{pmatrix}
		m_{SS}(1-\chi)^2\,\mathbf{C}_{SS} &
		m_{SI}\,\chi(1-\chi)\,\mathbf{C}_{SI} \\[4pt]
		m_{IS}\,\chi(1-\chi)\,\mathbf{C}_{IS} &
		m_{II}\,\chi^2\,\mathbf{C}_{II}
	\end{pmatrix}.
	\label{eq:C}
\end{equation}
where the binary \textit{block mask} $\mathbf{m}=(m_{SS},m_{SI},m_{IS},m_{II})\in\{0,1\}^4$,
$m_{SI}=m_{IS}$ is bibliometrically meaningful for
source-only $(1,0,0,0)$, institution-only $(0,0,0,1)$,
bipartite $(0,1,1,0)$, and full-joint $(1,1,1,1)$ combined blocks.
The mixing weight $\chi$ controls the source--institution balance
in the bipartite and full-joint cases and is inert in the single-mode cases. We could name $\mathbf{C}$ the reference, attention or citation matrix and prefer the latter. 

The model parameters are summarised in Table \ref{tab:params}. Note that the corpus excludes (a) census work references not among target works, (b) non-census works that reference target works, and (c) authorships that include an omitted author or institution. 

\begin{table}[h]
	\centering
	\caption{Model parameters $\boldsymbol{\theta}$.}
	\label{tab:params}
	\begin{tabularx}{\textwidth}{lllX}
		\toprule
		Symbol & Name & Type & Role \\
		\midrule
		\multicolumn{4}{@{}l@{}}{\textit{Corpus construction}} \\[2pt]
		$\mathfrak{S}$           & source set            & set             & define the field, subject, category \\
		$[t^c_0,t^c_1]$          & census window         & interval        & works whose references are traced \\
		$[t^t_0,t^t_1]$          & target window         & interval        & works eligible to receive references \\
		$\tau^U$                 & institution threshold & integer         & minimum work count to retain an institution \\
		$\tau^A$                 & author threshold      & integer         & minimum work count to retain an author \\[4pt]
		\multicolumn{4}{@{}l@{}}{\textit{Matrix construction}} \\[2pt]
		$\delta\in\{0,1\}$       & reference normalisation & binary        & 0: a work contributes $\bar{R}$ units; \newline 1: a reference contributes 1 unit \\
		$\mathbf{m}\in\{0,1\}^4$ & block mask            & binary 4-vector & select active unit-type pairs ($SS$, $SI$, $IS$, $II$) \\
		$\chi\in[0,1]$           & mixing weight         & continuous      & source--institution balance within each work \\[4pt]
		\multicolumn{4}{@{}l@{}}{\textit{Ranking}} \\[2pt]
		$\alpha\in[0,1)$         & discount factor       & continuous      & downweight long walks; ensure convergence\\
		\bottomrule
	\end{tabularx}
\end{table}

\paragraph{Spectral ranking.}

Let $\mathbf{r} = r_p=\sum_q C_{pq}$ be the total attention from unit $p$,
$\mathbf{D}_r=\mathrm{diag}(r_p)$ the corresponding diagonal matrix, and
$\mathbf{p}$ a boundary condition \citep{vignaSpectralRanking2016} with entries
$p_p=a_p/\sum_q a_q$,
where $a_p=N^s_p$ for sources and $a_p=M_k$ for institutions.
The row-stochastic citation matrix is
$\mathbf{H}=\mathbf{D}_r^{-1}\mathbf{C}$.

The spectral ranking $\boldsymbol{\pi}$ \citep[see ][for an overview of the concept]{vignaSpectralRanking2016} is the largest eigenvector of the problem
\begin{equation}
	\boldsymbol{\pi}
	\;=\;
	\alpha\mathbf{H}^\top\boldsymbol{\pi}
	+(1-\alpha)\mathbf{p},
	\qquad
	\mathbf{1}^\top\boldsymbol{\pi}=1,
	\quad
	\boldsymbol{\pi}\ge 0,
	\label{eq:spectral}
\end{equation}
with damping factor $\alpha\in[0,1)$.
In the limit $\alpha\to 1$, equation~\eqref{eq:spectral} reduces to the
eigenvector problem
\begin{equation}
	\boldsymbol{\pi}^{*} = \mathbf{H}^\top\boldsymbol{\pi}^{*},
	\qquad
	\mathbf{1}^\top\boldsymbol{\pi}^{*}=1,
	\quad
	\boldsymbol{\pi}^{*}\ge 0,
	\label{eq:eig}
\end{equation}
whose solution is unique when $\mathbf{H}$ is primitive.
In that limit, \cite{gellerCitationInfluenceMethodology1978} notes that $\boldsymbol{\pi}^{*}$ is the limiting stationary distribution of a irreducible, aperiodic Markov chain with transition matrix $\mathbf{H}$. She shows that $\boldsymbol{\pi}^{*}$ is the \textit{influence weight} defined by \cite{pinskiCitationInfluenceJournal1976}, scaled by the reference count proportion $\mathbf{r}/|\mathbf{r}|$. 

The spectral ranking $\boldsymbol{\pi}$ depends on the number of works in each unit, while 
\begin{equation}
	\mathbf{v}
	\;=\;
	\mathbf{A}^{-1}\boldsymbol{\pi},
	\qquad
	\mathbf{A}=\mathrm{diag}(p_r)
	\label{eq:vstar}
\end{equation}
removes this dependence. \cite{palacios-huertaMeasurementIntellectualInfluence2004} show that $\mathbf{v}$ is a ranking that is uniquely invariant to several reasonable constraints \cite[but see][]{serranoMeasurementIntellectualInfluence2004}.

\paragraph{Computation via power iteration.}
To compute $\boldsymbol{\pi}$ we define the modified matrix
\begin{equation}
	\mathbf{\tilde{H}}
	\;=\;
	\alpha\mathbf{H}+(1-\alpha)\mathbf{p}\mathbf{1}^\top,
	\label{eq:Hstar}
\end{equation}
so that the fixed-point condition
$\boldsymbol{\pi}=\mathbf{\tilde{H}}^{\top}\boldsymbol{\pi}$
is equivalent to equation~\eqref{eq:spectral}.
Because $\mathbf{p}>0$ and $0<\alpha<1$, every entry of $\mathbf{\tilde{H}}$
is strictly positive, ensuring it is a primitive row-stochastic
matrix. By the Perron--Frobenius theorem, $\mathbf{\tilde{H}}^{\top}$
therefore has a unique stationary distribution $\boldsymbol{\pi}$,
recovered by power iteration
\begin{equation}
	\boldsymbol{\pi}^{(k+1)}
	\;=\;
	\mathbf{\tilde{H}}^{\top}\boldsymbol{\pi}^{(k)}
	\label{eq:power}
\end{equation}
and converging to $\boldsymbol{\pi}$ from any strictly positive starting
vector $\boldsymbol{\pi}^{(0)}$.

\paragraph{Computation.}
$\mathbf{C}(\chi,\mathbf{m},\rho)$ is a sparse matrix built from a weighted edge list $\mathcal{E}$
accumulated over $\mathcal{R}$ with per-reference weight $\delta_i(\rho)$,
stored in CSR format, and row-normalised to $\mathbf{H}$.
The pipeline
$\mathcal{R}\!\to\!\mathcal{E}\!\to\!\mathbf{C}\!\to\!\mathbf{H}\!\to\!\boldsymbol{\pi}^*\!\to\!\mathbf{v}^*$ is implementation using \texttt{scipy.sparse.csr\_matrix} with parameter entry points detailed in the Supplement.
